\documentclass{article}
\usepackage{graphicx} % Required for inserting images
\usepackage[T1]{fontenc}
\usepackage{amsmath}
\usepackage[margin=1in]{geometry}
\usepackage{booktabs}
\usepackage{tabularx}
\usepackage[hidelinks]{hyperref}

\title{Lesson 9: Backtesting Without Fooling Yourself}
\author{Analyst onboarding \textperiodcentered{} Lesson 9 of 10 \textperiodcentered{} L/S Gamma Desk, QUANTT}
\date{2026-10-02}

\begin{document}

\maketitle

\textbf{Goal:} learn how to test a strategy on history honestly, and the traps that make most backtests look better than reality.

\textbf{You need:} Lessons 7--8 (forecasting, our rules).

\section{What a backtest is}

A \textbf{backtest} replays a strategy's rules on historical data, day by day, as if you had traded them, and records the P\&L. It is how we decide whether a rule (a band, a forecaster, an exit) is worth trading.

The danger: \textbf{a backtest is easy to fool, and the person most easily fooled is the one who wrote it.} Three traps matter most.

\section{Trap 1: look-ahead bias}

\textbf{Look-ahead bias} means the simulation uses information that wasn't available at the moment of the decision. The results look great and can never be repeated live.

Examples relevant to our desk:

\begin{itemize}

\item \textbf{Fitting GARCH on the full history} and then ``forecasting'' 2018 with parameters that already saw 2020.

\item \textbf{Using today's closing VIX at 10:15 am.} VIX closes at 4:15 pm, so at decision time only yesterday's value exists.

\item \textbf{Using today's partial price bar}, or filling trades at prices from \emph{before} the decision was made.

\item \textbf{Regime models that peek:} labelling day $t$ with a model fitted on later days.

\end{itemize}

\textbf{The rule:} for every number used on day $t$, ask \emph{``could I have known this at 10:15 am on day t?''} If not, it's look-ahead.

\section{Trap 2: overfitting}

\textbf{Overfitting} means tuning rules until they fit the past's noise rather than a real pattern. Try 50 band settings and one will look brilliant by luck alone.

\begin{itemize}

\item Studies of hundreds of trading algorithms found that backtest Sharpe ratios barely predict live results. The more backtests someone ran, the bigger the gap.

\item With 5 years of daily data, trying about 45 independent variations makes a Sharpe ratio of 1 or more \textbf{likely by chance alone} (Bailey et al.).

\end{itemize}

\textbf{The Sharpe ratio}, the standard score, is average return divided by volatility of returns, annualized. Above 1 is good. But a Sharpe ratio means little unless you know \textbf{how many things were tried} to find it.

\section{Trap 3: unrealistic trading costs}

Option spreads are wide. A backtest that trades at the mid with no costs is fantasy. We assume fills at \textbf{mid \ensuremath{\pm} slippage} (the same \$0.05 the live system uses) and include commissions on hedge trades.

\section{The fix: walk-forward testing plus a locked hold-out}

\textbf{1. Walk-forward optimization.} Never tune and test on the same data:

\begin{center}\small
\begin{tabularx}{\linewidth}{l >{\raggedright\arraybackslash}X >{\raggedright\arraybackslash}X l}
\toprule
\textbf{Fold} & \textbf{Tune rules on (train)} & \textbf{Gap (embargo)} & \textbf{Test on} \\
\midrule
1 & 2014--2017 & \textasciitilde{}30 trading days & 2018 \\
2 & 2014--2018 & \textasciitilde{}30 trading days & 2019 \\
3 & 2014--2019 & \textasciitilde{}30 trading days & 2020 \\
... & ... & ... & ... \\
\bottomrule
\end{tabularx}
\end{center}

Each test year uses settings chosen \emph{only} from earlier years. Stitch the test years together for an honest result.

\textbf{2. Embargo.} A trade lasts about 23 days, so a trade opened at the end of training would still be open in the test year. The gap stops results leaking across the boundary.

\textbf{3. Final hold-out.} Lock away the last \textasciitilde{}12 months. Run the final, chosen strategy on it \textbf{exactly once}. If you look, tweak, and look again, it is no longer out-of-sample.

\textbf{4. Log every trial.} Record every configuration tried, not just the winner. The \textbf{Deflated Sharpe Ratio} then corrects the best Sharpe for how many trials it took to find it.

\section{Our data, and its limits}

\begin{itemize}

\item \textbf{The club's R2 bucket:} intraday SPY option \textbf{trades} from 2014 to 2026, and daily SPY prices (CRSP).

\item \textbf{No quotes and no intraday SPY price.} Daily option prices must be built from trades, and the bid-ask spread has to be modelled.

\item \textbf{Few independent trades:} about 12 years of roughly 23-day trades gives \textbf{only \textasciitilde{}130 non-overlapping trades}. Keep the number of tuned parameters tiny.

\end{itemize}

\section{How our backtester will be built}

\begin{itemize}

\item \textbf{The same code as live:} each simulated day calls the same \texttt{signals}, \texttt{forecasting}, \texttt{selection}, \texttt{exits} and \texttt{pnl} modules as the live runner, with a simulated broker. The backtest then tests the exact rules we trade.

\item \textbf{Point-in-time inputs only:} the forecast is refit on returns up to the previous close, and VIX is lagged one day.

\item \textbf{Every run logged automatically} to a trial registry: settings, code version and daily returns.

\end{itemize}

\section{What to report}

Report Sharpe, plus \textbf{Sortino} (counts only downside volatility), \textbf{maximum drawdown}, \textbf{skew} (short-vol strategies have rare big losses), cost drag, and the \textbf{number of trials} behind the result. Show the walk-forward results and the hold-out result separately.

\section{Further reading}

Our report on this exact topic is \texttt{research/\allowbreak{}strategy\_\allowbreak{}research/\allowbreak{}2026-\allowbreak{}10-\allowbreak{}01-\allowbreak{}walk-\allowbreak{}forward-\allowbreak{}backtest-\allowbreak{}without-\allowbreak{}look-\allowbreak{}ahead.\allowbreak{}pdf}.

\section{Words to know}

\begin{itemize}

\item \textbf{Look-ahead bias:} using information not available at decision time.

\item \textbf{Overfitting:} fitting noise, not signal.

\item \textbf{Walk-forward:} tune on the past, test on the next unseen period, roll forward.

\item \textbf{Embargo:} gap between train and test periods. \textbf{Hold-out:} final data used only once.

\item \textbf{Sharpe /\allowbreak{} Sortino /\allowbreak{} drawdown:} return per unit of risk /\allowbreak{} per unit of downside risk /\allowbreak{} fall from a peak.

\end{itemize}

\section{Check yourself}

\begin{enumerate}

\item Your backtest decides at 10:15 am using that day's GARCH fit on closing prices. What's wrong?

\item You test 40 band settings and report the best Sharpe of 1.4. What else must you report?

\item Why must the hold-out be used only once?

\end{enumerate}

\emph{Answers: (1) Look-ahead: the close isn't known at 10:15. Fit on data up to the previous close. (2) That 40 were tried, and a Deflated Sharpe. (3) Every look lets you adapt to it, so it stops being unseen data.}

\end{document}
